\section{Integrated Gradients trên một đường đi}
\label{sec:ch05-integrated-gradients}

Hai tiên đề dẫn tới Integrated Gradients. Cho $x_0$ là baseline và $x$ là đầu
vào cần giải thích. Phương pháp chọn đoạn thẳng $x_0+\alpha(x-x_0)$ khi
$\alpha$ chạy từ không tới một, lấy gradient tại mọi điểm trên đoạn ấy, rồi
tích theo từng tọa độ:

\begin{equation}
\label{eq:ch05-integrated-gradients}
\operatorname{IG}_i(x;x_0)=(x_i-x_{0i})
\int_0^1 \frac{\partial F\bigl(x_0+\alpha(x-x_0)\bigr)}{\partial x_i}
\,d\alpha.
\end{equation}

Thừa số $x_i-x_{0i}$ biến độ dốc thành phần thay đổi theo tọa độ. Vì thế,
phương trình~\ref{eq:ch05-integrated-gradients} không đọc riêng độ dốc ở $x$;
nó cộng các độ dốc gặp trên cả hành trình từ $x_0$ đến $x$. Hình~\ref{fig:ch05-hai-phep-tich}
đặt hành trình đó cạnh phép tổng hợp không gian mà Grad-CAM sẽ dùng.

\begin{proposition}
\label{prop:ch05-tinh-day-du}
Nếu $F$ liên tục trên đoạn từ $x_0$ đến $x$, khả vi hầu khắp nơi và các đạo
hàm riêng của nó khả tích trên đoạn, tổng các attribution Integrated Gradients
bằng chênh lệch đầu ra:
\begin{equation}
\sum_i \operatorname{IG}_i(x;x_0)=F(x)-F(x_0).
\end{equation}
\end{proposition}

Mệnh đề~\ref{prop:ch05-tinh-day-du} là \tn{tính đầy đủ}{completeness}. Nó là
định lý cơ bản của giải tích áp cho đường đi trong không gian đầu vào. Khi chỉ
một tọa độ thay đổi, tổng trên chỉ còn tọa độ ấy; do đó tính đầy đủ kéo theo \tn{độ
nhạy}{sensitivity} của mục~\ref{sec:ch05-hai-tien-de}. Mệnh đề không nói từng phần đóng góp
là duy nhất. Những đường đi khác giữa cùng hai điểm cũng có thể thỏa tính đầy
đủ, nhưng chia chênh lệch cho các tọa độ theo cách khác.

Trong máy tính, tích phân được xấp xỉ bằng tổng gradient tại nhiều điểm trên
đoạn thẳng. Sai số xấp xỉ phải được kiểm tra bằng chênh lệch giữa tổng
attribution và $F(x)-F(x_0)$, thay vì mặc định rằng một số bước bất kỳ là đủ.
Đó là một kiểm tra số học của phép tính, không phải xác nhận rằng baseline đã
đại diện cho sự vắng mặt của tín hiệu.

\index{Integrated Gradients|(}%
\index{tích phân đường}%
\index{tính đầy đủ}%
